\documentclass{article}
\input{~/studies/preamble}

\title{Clustering}
\author{Hyoung Joon, Kim}
\date{\today}

\begin{document}

\maketitle
\tableofcontents
\newpage

\section{Introduction}
    We discuss about clustering, which is an unsupervised learning technique which identifies groups or clusters
    of data points.

\section{K-means Clustering}
    \subsection{Introduction}
        \textbf{K-means Clustering} is one of the most simple clustering algorithm. It does not use any probabilistic notions,
        and only rely on deterministic informations. K-means clustering can be seen as special case of \emph{Gaussian mixture model},
        , a probabilistic clustering algorithm, which will be shown later.
    
    \subsection{Algorithm}
        \begin{algorithm}
            \DontPrintSemicolon
            \SetArgSty{textnormal}
            \SetKwInOut{Input}{Input}
            \SetKwInOut{Output}{Output}
            
            \caption{K-means clustring algorithm}

            \Input{$\{\textbf{x}_1,\dots,\textbf{x}_N\}$ \newline
            number of clusters $K$ \newline
            number of iterations $T$}
            \Output{cluster vectors $\bm\mu_k$ \newline
            indicator \{$r_{nk}$\}}

            \BlankLine
            \hrule
            \tcp{initialize}
            $\{r_{nk} \leftarrow 0\}$ for $n \in \{1,\dots,N\}$ and $k \in \{1,\dots,K\}$\;
            \Repeat{$\{r_{nk}\}=\{r_{nk}^{(\text{old})}\}$ or iterated $T$ times}
            {
                \tcp{E step}
                $\{r_{nk}^{(\text{old})}\}$\;
                \For{$n \in \{1,\dots,N\}$}
                {
                    $k \leftarrow \arg\min_j\|\textbf{x}_n - \bm\mu_j\|^2$\;
                    $r_{nk} \leftarrow 1$\;
                    $\{r_{nk} \leftarrow 0\}$ for $j \in \{1,\dots,K\},\ j \neq k$
                }
                \tcp{M step}
                \For{$k \in \{1,\dots,K\}$}
                {
                    $\bm\mu_k \leftarrow \frac{\sum_n r_{nk}\textbf{x}_n}{\sum_n r_{nk}}$\;
                }
            }
            \Return{$\{\bm\mu_k\}$, $\{r_{nk}\}$}

        \end{algorithm}
        
        The algorithm tries to minimize the error function $\displaystyle{J=\sum_{n=1}^{N}\sum_{k=1}^{K}r_{nk}\|\textbf{x}_n-\bm\mu_k\|^2}$,
        where $r_{nk}$ is 1 if data point $\textbf{x}_n$ is included in cluster $k$, else 0.
        and $\bm\mu_k$ are the 'prototype' of the cluster vector, representing the center of the clusters. The algorithm first minimize the loss function
        about $r_{nk}$ with $\bm\mu\_k$ fixed, which is simply
        \[r_{nk}=
        \begin{cases}
            1 & \text{if }k=\arg\min_j \|\textbf{x}_n-\bm\mu_j\|^2 \\
            0 & \text{otherwise}
        \end{cases}\]
        This is \textbf{E step} if we look it as a kind of EM algorithm. Then, the algorithm fixes $r_{nk}$ and minimize about
        $\bm\mu_k$. The objective function $J$ is a quadratic function of $\bm\mu_k$ and can be minimized easily, with
        \[
        \bm\mu_k = \frac{\sum_n r_{nk}\textbf{x}_n}{\sum_n r_{nk}}.
        \]

        We can derive a sequential version of this batchwise algorithm.
        
        \begin{algorithm}
            \DontPrintSemicolon
            \SetArgSty{textnormal}
            \SetKwInOut{Input}{Input}
            \SetKwInOut{Output}{Output}
            
            \caption{Sequential K-means clustring algorithm}

            \Input{$\{\textbf{x}_1,\dots,\textbf{x}_N\}$ \newline
            number of clusters $K$ \newline
            number of iterations $T$}
            \Output{cluster vectors $\bm\mu_k$ \newline
            indicator \{$r_{n}$\}}

            \BlankLine
            \hrule
            \tcp{initialize}
            $\{r_{n} \leftarrow 0\}$ for $n \in \{1,\dots,N\}$ and $k \in \{1,\dots,K\}$\;
            $\{N_k \leftarrow 0\}$\;
            \ForEach{$\textbf{x} \in \{\textbf{x}_1,\dots,\textbf{x}_N\}$}
                {
                    $k \leftarrow \arg\min_j\|\textbf{x}_n - \bm\mu_j\|^2$\;
                    $N_k \leftarrow N_k + 1$\;
                    $r_n \leftarrow k$\;
                    $\bm\mu_k \leftarrow \bm\mu_k + \frac{1}{N_k}(\textbf{x}_n-\bm\mu_k)$\;
                }
            \Return{$\{\bm\mu_k\}$, $\{r_{n}\}$}

        \end{algorithm}
        The sequential k-means algorithm does not change the data points which are already clustered,
        unlike the original k-means algorithm, and the indicator changed from one-hot encoding to integer value.
    \subsection{Applications}
        K-means clustering can be used at image segmentation, specifically color segmentation.
\section{Mixtures of Gaussians}
    The \textbf{Gaussian mixture model} (GMM) is a simple linear superposition of Gaussian components. This can be formulated in
    terms of discrete latent variables. See this \href{run:../../03_bayesian_ml/models/gmm/gmm.pdf}{document} for more information.
\end{document}